% القسم: المناهج
% الفصل: البراهين
% المبحث: قراءة البراهين

\documentclass[../../../include/open-logic-section]{subfiles}

\begin{document}

\olfileid{mth}{prf}{rea}
\olsection{قراءة البراهين}

نادراً جداً ما تذكر البراهين التي تجدها في الكتب والمقالات جميع
التفاصيل التي أدرجناها حتى الآن في أمثلتنا. فكثيراً ما لا ينبه المؤلفون
إلى مواضع تمييز الحالات أو إقامة البرهان غير المباشر، ولا يذكرون أنهم
استعملوا تعريفاً. ولذلك ستحتاج غالباً، عندما تقرأ برهاناً في كتاب، إلى
ملء تلك التفاصيل بنفسك لكي تفهمه. وهذا تدريب جيد أيضاً على إتقان مختلف
الخطوات التي تحتاج إلى اتخاذها في البرهان. فلننظر في مثال.

\begin{prop}[الامتصاص]
لكل مجموعتين $A$ و$B$،
\[
A \cap (A \cup B) = A.
\]
\end{prop}

\begin{proof}
إذا كان $z \in A \cap (A \cup B)$، فإن $z \in A$، ومن ثم
$A \cap (A \cup B) \subseteq A$. والآن افترض $z \in A$. عندئذ يكون
$z \in A \cup B$ أيضاً، ولذلك يكون $z \in A \cap (A \cup B)$ أيضاً.
\end{proof}

البرهان السابق لقانون الامتصاص شديد الاختصار. فلا يذكر أياً من التعريفات
المستعملة، ولا يقول «علينا أن نثبت أن» قبل إثبات شيء، وهكذا. فلنفككه.
القضية المثبتة قضية كلية عن أي مجموعتين $A$ و$B$، وعندما يذكر البرهان
$A$ أو $B$ فهما متغيران لمجموعتين اعتباطيتين. والقضية الكلية التي
يثبتها البرهان هي ما يلزم لإثبات تطابق مجموعتين: أن كل !!{element} في
الطرف الأيسر من المتطابقة هي !!a{element} في الطرف الأيمن، والعكس.

\begin{quote}
«إذا كان $z \in A \cap (A \cup B)$، فإن $z \in A$، ومن ثم
$A \cap (A \cup B) \subseteq A$».
\end{quote}

هذا هو النصف الأول من برهان المتطابقة: فهو يثبت أنه إذا كان $z$
الاعتباطي !!a{element} في الطرف الأيسر، فإنه أيضاً !!a{element} في
الطرف الأيمن؛ أي $A \cap (A \cup B) \subseteq A$. افترض أن
$z \in A \cap (A \cup B)$. وبما أن $z$ يكون !!{element} في تقاطع
مجموعتين إذا وفقط إذا كان !!{element} في كلتيهما، نستنتج أن
$z \in A$ وكذلك $z \in A \cup B$. وبوجه خاص $z \in A$، وهذا ما أردنا
بيانه. ولما كان ذلك كل ما ينبغي فعله في النصف الأول، نعلم أن بقية البرهان
لا بد أن تكون برهاناً للنصف الثاني، أي برهاناً على
$A \subseteq A \cap (A \cup B)$.

\begin{quote}
«والآن افترض $z \in A$. عندئذ يكون $z \in A \cup B$ أيضاً، ولذلك يكون
$z \in A \cap (A \cup B)$ أيضاً».
\end{quote}

نبدأ بافتراض $z \in A$، لأننا نبين أنه لكل~$z$، إذا كان $z \in A$
فإن $z \in A \cap (A \cup B)$. ولكي نبين
$z \in A \cap (A \cup B)$، علينا أن نبين، بحسب تعريف «$\cap$»، أن
(1) $z \in A$، وأن (2) $z \in A \cup B$. أما (1) فهو فرضنا نفسه، فلا
شيء آخر ينبغي إثباته، ولذلك لا يذكره البرهان مرة أخرى. وبالنسبة إلى (2)،
تذكر أن $z$ يكون !!a{element} في اتحاد مجموعتين إذا وفقط إذا كان
!!{element} في واحدة منهما على الأقل. وبما أن $z \in A$ وأن
$A \cup B$ اتحاد $A$ و$B$، فهذا حاصل هنا. إذن $z \in A \cup B$.
لقد بينا (1) $z \in A$ و(2) $z \in A \cup B$، ومن ثم، بحسب تعريف
«$\cap$»، $z \in A \cap (A \cup B)$. لا يذكر البرهان تلك التعريفات؛
إذ يفترض أن القارئ قد استبطنها. فإذا لم تكن قد فعلت، فعليك الرجوع
وتذكير نفسك بها. وعليك كذلك أن تدرك لماذا ينتج من $z \in A$ أن
$z \in A \cup B$، ولماذا ينتج من $z \in A$ و$z \in A \cup B$ أن
$z \in A \cap (A \cup B)$.

وفيما يأتي صيغة أخرى للبرهان السابق تصرح بكل شيء:
\begin{proof}{}
[بحسب تعريف $=$ للمجموعات، لإثبات $A \cap (A \cup B) = A$ علينا أن
  نبين (أ) $A \cap (A \cup B) \subseteq A$ و(ب)
  $A \subseteq A \cap (A \cup B)$. (أ): بحسب تعريف $\subseteq$، علينا
  أن نبين أنه إذا كان $z \in A \cap (A \cup B)$ فإن $z \in A$.]
إذا كان $z \in A \cap (A \cup B)$، فإن $z \in A$ [لأنه بحسب تعريف
$\cap$، يكون $z \in A \cap (A \cup B)$ إذا وفقط إذا كان $z \in A$
و$z \in A \cup B$]، ومن ثم $A \cap (A \cup B) \subseteq A$.
[(ب): بحسب تعريف $\subseteq$، علينا أن نبين أنه إذا كان $z \in A$ فإن
$z \in A \cap (A \cup B)$.] افترض الآن [(1)] $z \in A$. عندئذ أيضاً
[(2)] $z \in A \cup B$ [لأن (1) يعطينا $z \in A$ أو $z \in B$،
وهو ما يعني، بحسب تعريف $\cup$، أن $z \in A \cup B$]؛ ومن ثم يكون
$z \in A \cap (A \cup B)$ أيضاً [لأن تعريف $\cap$ يقتضي
$z \in A$، أي (1)، و$z \in A \cup B$، أي (2)].
\end{proof}

\begin{prob}
فصّل البرهان الآتي لـ$A \cup (A \cap B) = A$، ذاكراً جميع أنماط
الاستدلال المستعملة، وسبب اتباع كل خطوة من الافتراضات أو القضايا المثبتة
قبلها، والمواضع التي نحتاج فيها إلى أي تعريف.
\begin{proof}
  إذا كان $z \in A \cup (A \cap B)$ فإن $z \in A$ أو
  $z \in A \cap B$. وإذا كان $z \in A \cap B$ فإن $z \in A$. وكل
  $z \in A$ ينتمي أيضاً إلى $A \cup (A \cap B)$.
\end{proof}
\end{prob}

\end{document}
